\section{Full support and regularity}
\label{sec:variation}

We now turn the constant-interval curvature into full support.
The first variation must allow removal of arbitrary Stieltjes mass,
not only removal of atoms: a minimizer need not initially be known
to have a discrete or absolutely continuous measure.
For related variational characterizations see
Chen, Handschy, and Lerman~\cite[Propositions~3--4]{CHL2018} and
Jagannath and Tobasco~\cite[Lemma~3.2 and Corollary~3.6]{JT2017}.

\begin{lemma}[First variation in every admissible integrable direction]
\label{lem:first-variation}
Let $\gamma\in\mathcal U$, and suppose that $\eta\in L^1(0,1)$
has the property that $\gamma+\epsilon\eta\in\mathcal U$ for
$0\leq\epsilon\leq\epsilon_0$, where $\epsilon_0>0$.
With $X$ and $q$ from Lemma~\ref{lem:diffusion},
\begin{equation}
 \left.\frac{\dd}{\dd\epsilon}\right|_{\epsilon=0+}
       \mathcal P(\gamma+\epsilon\eta)
   =\frac12\int_0^1\eta(t)(q(t)-t)\,\dd t.
 \label{eq:first-variation}
\end{equation}
\end{lemma}

\begin{proof}
Put $\beta=\gamma+\epsilon\eta$, $u=\Phi_{\gamma,x}$,
$\widetilde u=\Phi_{\beta,x}$, and
$D=\Phi_\beta-\Phi_\gamma$.  Subtracting the two equations gives
\[
 D_t=-\tfrac12D_{xx}
       -\tfrac12\beta(\widetilde u+u)D_x
       -\tfrac12(\beta-\gamma)u^2.
\]
Let $Y^\epsilon_0=0$ and
\[
 \dd Y^\epsilon_t
   =\tfrac12\beta(t)(\widetilde u+u)(t,Y^\epsilon_t)\,\dd t
      +\dd B_t.
\]
This diffusion exists uniquely by Lemma~\ref{lem:regularity};
its drift is bounded by the integrable function $\beta$.
It\^o's formula on $[0,T]$ yields the following identity.  One
may justify it without a classical time derivative by approximating
$\gamma$ and $\beta$ separately by nondecreasing step functions,
using the same smooth terminal datum for the pair.  The two
solutions, their spatial derivatives, and the associated linear
diffusions converge on $[0,T]$ by Lemma~\ref{lem:regularity} and
the same Gronwall estimate as in Lemma~\ref{lem:diffusion}.
The source converges in $L^1$, since the gradients are bounded
by one.  Passing to the limit in their classical identities gives
\[
 D(0,0)=\E D(T,Y^\epsilon_T)
       +\frac12\int_0^T(\beta-\gamma)(t)
                       \E[u(t,Y^\epsilon_t)^2]\,\dd t.
\]
The time-local form of \eqref{eq:value-stability} gives
$\|D(T,\cdot)\|_\infty
 \leq\frac32\int_T^1|\beta-\gamma|$.
Letting $T\uparrow1$ therefore proves the exact identity
\begin{equation}
 \frac{\Phi_\beta(0,0)-\Phi_\gamma(0,0)}\epsilon
    =\frac12\int_0^1\eta(t)
                     \E[u(t,Y^\epsilon_t)^2]\,\dd t.
 \label{eq:exact-difference}
\end{equation}

The coefficients $\beta$ have uniformly bounded $L^1$ norms for
$0\leq\epsilon\leq\epsilon_0$, and converge to $\gamma$ in
$L^1$.  Lemma~\ref{lem:regularity} gives uniform convergence
$\widetilde u\to u$ and common spatial Lipschitz bounds on every
compact strip.  Coupling the diffusions with the same Brownian
motion and using Gronwall's inequality gives
$Y^\epsilon\to X$ uniformly there.  The integrand in
\eqref{eq:exact-difference} consequently converges for each $t<1$
to $\eta(t)q(t)$ and is bounded in absolute value by $|\eta(t)|$.
Dominated convergence and differentiation of the linear penalty
term in $\mathcal P$ prove \eqref{eq:first-variation}.
\end{proof}

\begin{lemma}[Contact conditions]
\label{lem:contact}
Let $\gamma$ minimize $\mathcal P$ over $\mathcal U$, and let
$\mu$ be its locally finite Stieltjes measure on $[0,1)$, with the
convention $\mu(\{0\})=\gamma(0)$.  Thus
$\gamma(t)=\mu([0,t])$.  Define
\begin{equation}
 G(s)=\int_s^1(q(t)-t)\,\dd t,
 \qquad S=\operatorname{supp}\mu\subset[0,1).
 \label{eq:contact-potential}
\end{equation}
Then $G\geq0$ on $[0,1)$ and $G=0$ on $S$.  At every
$s\in S\cap(0,1)$,
\begin{equation}
 q(s)=s,\qquad q'(s)\leq1.
 \label{eq:contact-derivatives}
\end{equation}
If $0\in S$, then also $q'(0)\leq1$, with the derivative from
the right.
\end{lemma}

\begin{proof}
For $s<1$, the direction $\eta=\mathbf1_{[s,1)}$ adds an atom
to $\mu$ and is admissible.  Minimality and
Lemma~\ref{lem:first-variation} give $G(s)\geq0$.
If $K\subset[0,1)$ is compact, let $\nu=\mu|_K$ and take
$\eta_K(t)=-\nu([0,t])$.  This bounded direction is admissible
for $0\leq\epsilon\leq1$, since it replaces $\mu$ by
$\mu-\epsilon\nu$.  Fubini's theorem gives
\[
 0\leq D\mathcal P(\gamma)[\eta_K]
      =-\frac12\int_K G(s)\,\mu(\dd s).
\]
The interchange is legitimate because $|q(t)-t|\leq1$ and
$\nu(K)<\infty$.  Since $G\geq0$, it vanishes
$\mu$-almost everywhere on every such $K$.  Its continuity
then gives $G=0$ on the relative support $S$.

By Lemma~\ref{lem:diffusion},
$G'(s)=s-q(s)$ and $G''(s)=1-q'(s)$ on compact strips.
Every interior support point is a minimum of $G$, proving
\eqref{eq:contact-derivatives}.  If $0\in S$, use $G(0)=0$
and $q(0)=0$ to write
\[
 G(h)=\tfrac12(1-q'(0))h^2+o(h^2).
\]
Nonnegativity for $h>0$ proves the asserted one-sided condition.
\end{proof}

\begin{lemma}[No terminal plateau]
\label{lem:no-terminal-plateau}
For a minimizer $\gamma$, there is no $a<1$ and finite $c\geq0$
such that $\gamma=c$ on $(a,1)$.
\end{lemma}

\begin{proof}
The endpoint conclusion is recorded by Auffinger, Chen, and
Zeng~\cite[Remark~7]{ACZ2017}. We give the quantitative argument
needed here, including arbitrary positive constant plateaus.
The boundary-layer method is also related to the zero-terminal-segment
argument in El Alaoui, Montanari, and
Sellke~\cite[Lemma~6.12 and the proof of Theorem~5 in the cited preprint]{EAMS2021};
their positivity-support statement concerns a different admissible class.
Suppose such a plateau exists.  We show that for some $C>0$,
\begin{equation}
 1-q(t)\geq C\sqrt{1-t}
 \quad\hbox{for all $t$ sufficiently close to one}.
 \label{eq:terminal-layer}
\end{equation}
Fix $s\in(a,1)$, and let $\rho_s$ denote the law of $X_s$.
Choose $R<\infty$ with $p_0=\rho_s([-R,R])>0$; no density
assumption on $\rho_s$ is needed.

When $c>0$, set
$h(t,y)=\E\exp(c|y+\sqrt{1-t}Z|)$, with $Z$ standard normal.
The heat formula gives $\Phi(t,y)=c^{-1}\log h(t,y)$ and
drift $c\Phi_x=h_x/h$.  Its transition density on the plateau is
\begin{equation}
 k_{s,t}(x,y)=p_{t-s}(y-x)\frac{h(t,y)}{h(s,x)},
 \qquad
 p_v(z)=\frac{e^{-z^2/(2v)}}{\sqrt{2\pi v}}.
 \label{eq:terminal-kernel}
\end{equation}
Indeed, the heat-semigroup identity makes
$h(t,B_t)/h(s,B_s)$ a positive mean-one martingale when
Brownian motion is started at $(s,x)$.  Its stochastic logarithm
has the bounded coefficient $h_x/h=c\Phi_x$, so change of
measure gives precisely this drift and kernel.
For $c=0$ use the same kernel formula with $h\equiv1$.

Put $L=1-s$ and $\delta=L/2$.  Uniformly for
$s+\delta\leq t<1$, $|x|\leq R$, and $|y|\leq1$,
\[
 p_{t-s}(y-x)\geq
 (2\pi L)^{-1/2}e^{-(R+1)^2/(2\delta)}=:k_*>0,
\]
while $h(t,y)\geq1$ and
$h(s,x)\leq e^{cR}\E e^{c\sqrt L|Z|}=:M$.
Consequently the mixture of \eqref{eq:terminal-kernel} against
$\rho_s$ is a density $f_t$ satisfying
\begin{equation}
 f_t(y)\geq p_0k_*/M=:m>0
 \quad (s+\delta\leq t<1,\ |y|\leq1).
 \label{eq:terminal-density}
\end{equation}

For $\epsilon=1-t$ and $|y|\leq\sqrt\epsilon$,
differentiation of the terminal heat formula gives, also when
$c=0$,
\[
 u(t,y)=
 \frac{\E[\operatorname{sgn}(y+\sqrt\epsilon Z)
                      e^{c|y+\sqrt\epsilon Z|}]}
      {\E e^{c|y+\sqrt\epsilon Z|}}.
\]
For $\epsilon\leq1$ the denominator is at most
$D=e^c\E e^{c|Z|}<\infty$.  The unnormalized mass of each
sign is at least $p_*:=\Pr(Z\geq1)>0$, because
$|y/\sqrt\epsilon|\leq1$ and the exponential tilt is at least
one.  Each sign thus has probability at least
$\theta=p_*/D\in(0,1/2)$ under the normalized tilt.  Hence
$|u(t,y)|\leq1-2\theta$ and
$1-u(t,y)^2\geq4\theta(1-\theta)=:\kappa>0$.
Integrating this bound over
$[-\sqrt\epsilon,\sqrt\epsilon]$ and using
\eqref{eq:terminal-density} proves
\eqref{eq:terminal-layer} with $C=2m\kappa$.

It follows that
$q(t)-t\leq (1-t)-C\sqrt{1-t}<0$ sufficiently near one.
Thus $G(s')<0$ for $s'$ sufficiently near one, contrary to
Lemma~\ref{lem:contact}.  This also covers $c=0$.
\end{proof}

\begin{proof}[Proof of Theorem~\ref{thm:full-support}]
Let $S$ be the support in Lemma~\ref{lem:contact}.  Since $S$
is relatively closed in $[0,1)$, a nonempty component of its
complement is an initial, interior, or terminal interval.
On each such interval $\gamma$ is constant.  A terminal
component is impossible by Lemma~\ref{lem:no-terminal-plateau};
that lemma also excludes empty support, for which $\gamma=0$.

Consider a remaining gap with endpoints $0\leq a<b<1$ and
with $a,b\in S$.  By Lemma~\ref{lem:contact},
$q(a)=a$, $q(b)=b$, and $q'(a),q'(b)\leq1$, including the
one-sided interpretation at $a=0$.  If its constant coefficient
is positive, Proposition~\ref{prop:gap-curvature} makes $q'$
strictly convex inside the gap.  By continuity up to its endpoints,
\[
 q'(t)<\frac{b-t}{b-a}q'(a)+\frac{t-a}{b-a}q'(b)\leq1,
 \qquad a<t<b.
\]
Integrating contradicts $q(b)-q(a)=b-a$, as illustrated by the positive
area deficit in Figure~\ref{fig:gap-convexity}.
If the constant coefficient is zero, the strict increase in
Proposition~\ref{prop:gap-curvature} instead gives
$q'(t)<q'(b)\leq1$, with the same contradiction.

It remains only to consider an initial component $[0,b)$ when
$0\notin S$ and $b\in S$.  Here $\gamma=0$ on the gap.
Strict increase gives $q'(t)<q'(b)\leq1$ for $0<t<b$, whereas
$q(0)=0$ and $q(b)=b$.  Integration is again a contradiction.
Thus $S=[0,1)$.  The contact identities give $q(t)=t$ for
$0<t<1$, and the identity at zero follows from symmetry.
Differentiating this identity using Lemma~\ref{lem:diffusion} gives
$\E\Phi_{\gamma,xx}(t,X_t)^2=1$, including the right derivative
at zero. Finally, every nonempty interval $(a,b)$ with
$0\leq a<b<1$ has positive $\mu$ mass, so
$\gamma(b)-\gamma(a)=\mu((a,b])>0$.
\end{proof}
